\section{Множества и отношения}
\label{sec:sets-relations}

\subsection{Множества и операции над ними}

Множество --- совокупность различимых объектов. Запись $x\in A$ означает, что
$x$ является элементом $A$. Множество $A$ является подмножеством $B$, если
\[
    A\subseteq B
    \quad\Longleftrightarrow\quad
    \forall x\;(x\in A\Rightarrow x\in B).
\]

Основные операции:
\[
\begin{aligned}
A\cup B&=\{x:x\in A\text{ или }x\in B\},\\
A\cap B&=\{x:x\in A\text{ и }x\in B\},\\
A\setminus B&=\{x:x\in A,\ x\notin B\},\\
A\triangle B&=(A\setminus B)\cup(B\setminus A).
\end{aligned}
\]
Дополнение $\overline A$ рассматривается относительно фиксированного
универсального множества $U$.

Законы де Моргана:
\[
    \overline{A\cup B}=\overline A\cap\overline B,
    \qquad
    \overline{A\cap B}=\overline A\cup\overline B.
\]

Множество всех подмножеств $A$ обозначается $2^A$ или $\mathcal P(A)$.
Для конечного множества $|A|=n$ справедливо
\[
    |\mathcal P(A)|=2^n.
\]

\subsection{Мощность и декартово произведение}

Множества $A$ и $B$ равномощны, если существует биекция $A\to B$.
Множество счётно, если оно конечно или равномощно $\mathbb N$.
Множества $\mathbb N$, $\mathbb Z$, $\mathbb Q$ счётны, а $\mathbb R$ несчётно.

Декартово произведение определяется как
\[
    A\times B=\{(a,b):a\in A,\ b\in B\}.
\]
Для конечных множеств
\[
    |A\times B|=|A|\,|B|.
\]

\subsection{Бинарные отношения}

Бинарным отношением между $A$ и $B$ называется подмножество
\[
    R\subseteq A\times B.
\]
Если $(a,b)\in R$, пишут $aRb$. Область определения и множество значений:
\[
\operatorname{Dom}R=\{a\in A:\exists b\in B\;aRb\},
\qquad
\operatorname{Im}R=\{b\in B:\exists a\in A\;aRb\}.
\]

Обратное отношение:
\[
    R^{-1}=\{(b,a):(a,b)\in R\}.
\]
Композиция отношений $R\subseteq A\times B$ и $S\subseteq B\times C$:
\[
    S\circ R=
    \{(a,c):\exists b\in B\;(aRb\land bSc)\}.
\]
Композиция ассоциативна, но в общем случае некоммутативна.

\subsection{Свойства отношений на одном множестве}

Пусть $R\subseteq A\times A$. Отношение называется:
\begin{itemize}
    \item рефлексивным, если $aRa$ для всех $a\in A$;
    \item антирефлексивным, если $aRa$ не выполняется ни для одного $a$;
    \item симметричным, если $aRb\Rightarrow bRa$;
    \item антисимметричным, если $aRb$ и $bRa$ влекут $a=b$;
    \item транзитивным, если $aRb$ и $bRc$ влекут $aRc$;
    \item связным, если для любых различных $a,b$ выполнено $aRb$ или $bRa$.
\end{itemize}

Транзитивное замыкание $R^+$ --- наименьшее транзитивное отношение,
содержащее $R$:
\[
    R^+=R\cup R^2\cup R^3\cup\cdots.
\]
Рефлексивно-транзитивное замыкание обозначают $R^*$.

\subsection{Отношение эквивалентности}

Отношение эквивалентности рефлексивно, симметрично и транзитивно. Класс
эквивалентности элемента $a$:
\[
    [a]=\{x\in A:xRa\}.
\]
Классы эквивалентности либо совпадают, либо не пересекаются и образуют разбиение
множества $A$. Обратно, любое разбиение множества порождает отношение
эквивалентности: два элемента эквивалентны, если принадлежат одному блоку.

Множество всех классов обозначается $A/R$ и называется фактор-множеством.

\subsection{Отношения порядка}

Рефлексивное, антисимметричное и транзитивное отношение называется частичным
порядком. Пара $(A,\preceq)$ --- частично упорядоченное множество.
Если любые два элемента сравнимы, порядок линейный.

Следует различать:
\begin{itemize}
    \item минимальный элемент --- элемент, ниже которого нет другого;
    \item наименьший элемент --- элемент, не превосходящий всех остальных;
    \item максимальный и наибольший элементы;
    \item нижнюю и верхнюю границы подмножества;
    \item инфимум и супремум --- наибольшую нижнюю и наименьшую верхнюю границы.
\end{itemize}
Наименьший элемент, если существует, единственен; минимальных элементов может
быть несколько.

Конечный частичный порядок удобно изображать диаграммой Хассе, в которой
транзитивные рёбра и петли не показываются.

\subsection{Функции как специальные отношения}

Функция $f\colon A\to B$ --- отношение, в котором каждому $a\in A$ соответствует
ровно один $b\in B$. Функция называется инъективной, если разные аргументы имеют
разные значения; сюръективной, если каждый элемент $B$ является значением;
биективной, если выполняются оба свойства.

Обратная функция существует тогда и только тогда, когда исходная функция
биективна.

\subsection{Что важно сказать на экзамене}

Нужно показать связь трёх конструкций: операции над множествами, отношения как
подмножества декартова произведения и функции как частный случай отношений.
Особенно важны эквивалентность и разбиения, а также различие частичного и
линейного порядка.
